% Supplement to Evans 4.3: Radon and Laplace transforms, printed pp. 186--195.
\Needspace{35mm}
\subsection{Radon transform and plane-wave decomposition}\label{sec:4-radon}
For $f\in C_c^\infty(\R^n)$, set
\[
 Rf(s,\omega)=\int_{x\cdot\omega=s}f(x)\,dS_x,\qquad \omega\in S^{n-1}.
\]
Coordinates $x=s\omega+\sum y_jb_j$, where $b_j$ are an orthonormal
basis of $\omega^\perp$, make the surface measure $dy$.
They immediately give
\[
 Rf(-s,-\omega)=Rf(s,\omega),\quad
 R(\partial_jf)=\omega_j\partial_sRf,\quad R(\Delta f)=\partial_s^2Rf .
\]
For the middle identity decompose $\partial_j$ into a normal and
tangential derivative. Tangential integrals vanish by compact support.
Induction gives $R(D^\gamma f)=\omega^\gamma\partial_s^{|\gamma|}Rf$.
If $\operatorname{supp}f\subset B_R$, then $Rf=0$ for $|s|\geq R$.


\PDEBodyDiagram{4.3}{Radon data integrate along level hyperplanes.}
{In dimension two the integration sets are lines perpendicular to $\omega$. The scalar $s$ is signed distance in the normal direction, not distance traveled along a line. The blue line is one member of the family; higher dimensions replace it by an $(n-1)$-plane.}
{ch04-radon-lines}

\begin{theorem}[Slice theorem and inversion; source Theorems 4--5]
With the unnormalized one-dimensional transform
$B(r,\omega)=\int_{\R} e^{-irs}Rf(s,\omega)\,ds$,
\[
 B(r,\omega)=(2\pi)^{n/2}\widehat f(r\omega).
\]
For every dimension,
\[
 f(x)=\frac1{2(2\pi)^n}
 \int_{S^{n-1}}\int_{\R} B(r,\omega)|r|^{n-1}e^{irx\cdot\omega}\,dr\,dS_\omega .
\]
If $n=2k+1$ is odd, this becomes the local formula
\[
 f(x)=C_n\int_{S^{n-1}}
       \partial_s^{2k}Rf(x\cdot\omega,\omega)\,dS_\omega,\qquad
 C_n=\frac{(-1)^k}{2(2\pi)^{2k}} .
 \tag{RI4}
\]
\end{theorem}
\begin{expandedproof}
Fubini in the orthonormal coordinates gives
\[
 B(r,\omega)=\int f(x)e^{-ir\omega\cdot x}\,dx .
\]
In Fourier inversion, polar coordinates with signed $r$ double-count
each nonzero frequency, giving the factor $1/2$ and $|r|^{n-1}$.
All integrals converge since $\widehat f$ is rapidly decreasing.
For $n-1=2k$, one-dimensional inversion gives
\[
 \partial_s^{2k}Rf(s,\omega)
 =\frac{(-1)^k}{2\pi}\int_{\R} r^{2k}B(r,\omega)e^{irs}\,dr .
\]
Substituting this in the general inversion formula proves (RI4).
\end{expandedproof}
\Needspace{45mm}
\PDEAnnotation{奇偶维差别藏在绝对值里}{带符号的极径在一般维数必须用 $|r|^{n-1}$；只有奇维时才等于多项式 $r^{2k}$，从而变成有限阶局部导数。偶维的反演通常带非局部算子。}

For odd $n$, define $\rho_f=C_n\partial_s^{n-1}Rf$.
Then the plane-wave integral
\[
 u(x,t)=\int_{S^{n-1}}\rho_g(x\cdot\omega-t,\omega)\,dS_\omega
\]
solves $u_{tt}-\Delta u=0$, $u(0)=g$, $u_t(0)=0$.
The PDE holds because $|\omega|=1$; the first initial condition is
(RI4). The integrand for $u_t(0)$ is odd under $\omega\mapsto-\omega$,
since an odd $s$ derivative reverses the symmetry of $Rg$.

For the constant-coefficient system
$\mathbf u_t+\sum B_j\mathbf u_{x_j}=0$,
suppose $B(\omega)=\sum\omega_jB_j$ is diagonalizable with real
eigenvalues $\lambda_l(\omega)$.
In Radon coordinates the diagonal components satisfy one-dimensional
transport and equal
\[
 v_l(s,\omega,t)=h_l(s-\lambda_l(\omega)t,\omega).
\]
After applying (RI4), use
$\rho_l(s,\omega)=C_n\partial_s^{n-1}h_l(s,\omega)$ exactly once:
\[
 \mathbf u(x,t)=\int_{S^{n-1}}\sum_l
 a_l(\omega)\rho_l(x\cdot\omega-\lambda_l(\omega)t,\omega)\,dS_\omega .
\]
Here $a_l$ are the diagonalizing columns; the shift is not inserted
again in the definition of $\rho_l$.
For compact data in $B_R$, if
$\min_{\omega,l}|\lambda_l(\omega)|=\sigma>0$,
then $\mathbf u=0$ for $|x|<\sigma t-R$.
Indeed each shifted argument has magnitude at least $\sigma t-|x|>R$.
This is a conditional odd-dimensional Huygens statement; the
diagonalization and nonzero-speed hypotheses cannot be removed.
\Needspace{45mm}
\PDEAnnotation{平面支集条件怎样变成内部空区}{若 $|x|<\sigma t-R$，所有方向的平面参数都移出初始支集区间，整个反演积分为零。零特征速度会留下不移动的分量，因此必须单独排除。}

\Needspace{35mm}
\subsection{Laplace transform, resolvents and the heat--wave link}
If $v$ and $v_t$ have suitable exponential bounds, integration by parts gives
\[
 \mathcal L(v_t)(s)=s\mathcal Lv(s)-v(0),\qquad
 \mathcal Lv(s)=\int_0^\infty e^{-st}v(t)\,dt .
\]
Thus for $v_t-\Delta v=0$, $v(0)=f$,
\[
 (s-\Delta)\mathcal Lv=f .
\]
Use the same spatial boundary condition before and after transformation;
for example a zero Dirichlet condition stays zero.
In the whole space this recovers
$(s-\Delta)^{-1}f=\int_0^\infty e^{-st}(G_t*f)\,dt$.
For $L^2$ data the Bochner integral is justified by
$\|G_t*f\|_2\leq\|f\|_2$.
\Needspace{45mm}
\PDEAnnotation{初值正是 Laplace 分部积分的边界项}{时间导数并非只替换成乘以 $s$。零时刻留下 $-v(0)$，它成为椭圆 resolvent 方程的右端；漏掉这一项就把非零初值问题变成了齐次问题。}

\begin{theorem}[Gaussian averaging of the wave group; source Example 9]
Let $g$ be smooth compactly supported, and let $U(x,s)$ be the wave
solution with displacement $g$ and zero initial velocity, extended evenly
in $s$. Then
\[
 V(x,t)=\frac1{\sqrt{4\pi t}}\int_{\R} e^{-s^2/(4t)}U(x,s)\,ds
\]
is the heat solution with datum $g$.
\end{theorem}
\begin{expandedproof}
For smooth data the derivatives of $U$ grow at most polynomially,
as can also be checked from its Fourier representation.
Gaussian weights justify differentiation and kill the boundary terms.
Writing $q_t(s)=(4\pi t)^{-1/2}e^{-s^2/(4t)}$, observe
$\partial_tq_t=\partial_{ss}q_t$. Hence
\[
 V_t=\int q_{ss}U=\int qU_{ss}=\int q\Delta U=\Delta V.
\]
The Gaussian approximate identity in $s$ gives $V(0)=g$.
Alternatively the frequency multiplier is
$\int q_t(s)\cos(s|\xi|)\,ds=e^{-t|\xi|^2}$,
which proves the same statement in $L^2$.
\end{expandedproof}

For $n=2k+1\geq3$, let $M_g(x,r)$ be the spherical average of $g$.
Equating the two Gaussian representations gives, with
$\lambda=1/(4t)$,
\[
 \int_0^\infty e^{-\lambda s^2}U(x,s)\,ds
 =\frac{|S^{n-1}|}{2\pi^k}\lambda^k
       \int_0^\infty e^{-\lambda r^2}r^{2k}M_g(x,r)\,dr .
\]
Use $\lambda e^{-\lambda r^2}=-(2r)^{-1}\partial_re^{-\lambda r^2}$
and integrate by parts $k$ times. The result is
\[
 \lambda^k\int_0^\infty e^{-\lambda r^2}r^{2k}M_g\,dr
 =2^{-k}\int_0^\infty e^{-\lambda r^2}
       r(r^{-1}\partial_r)^k(r^{2k-1}M_g)\,dr .
\]
There are no boundary terms at infinity by compact support in $r$.
At zero, $M_g$ is smooth and even; before the last differentiation
the relevant odd power has not fallen below one, so the boundary
terms are zero.
After $s^2=\tau$, uniqueness of the Laplace transform implies
\[
 U(x,t)=\frac1{(n-2)!!}\,
 \partial_t(t^{-1}\partial_t)^{(n-3)/2}
                  (t^{n-2}M_g(x,t)).
\]
Here $|S^{n-1}|/(2^{k+1}\pi^k)=1/(n-2)!!$ follows from the
Gaussian sphere area formula.
The one-dimensional case is $U=(g(x+t)+g(x-t))/2$.
One can justify Laplace uniqueness without a new PDE theorem:
after multiplying by $e^{-a\tau}$, substitute $z=e^{-\tau}$;
vanishing at all $s=a+j$ gives zero moments on $[0,1]$.
Polynomial density in continuous functions makes the finite signed
measure zero. This applies to the continuous, exponentially integrable
functions above.
